\documentclass[12pt]{article}

% -----------------------------
% PAQUETES
% -----------------------------
\usepackage[spanish,es-noshorthands]{babel}
\usepackage[utf8]{inputenc}
\usepackage[T1]{fontenc}
\usepackage{lmodern}
\usepackage{amsmath,amssymb,amsthm}
\usepackage{enumitem}
\usepackage{geometry}
\usepackage{setspace}
\usepackage{xcolor}

% -----------------------------
% CONFIGURACIÓN
% -----------------------------
\geometry{margin=1.5cm}
\setstretch{1.15}

% -----------------------------
% DOCUMENTO
% -----------------------------
\begin{document}
	
	\begin{center}
		{\Large \textbf{Solución Problema 1}}\\[0.3cm]
		{\normalsize Pensamiento Matemático}\\[0.2cm]
		{\normalsize Contexto: Psicología}
	\end{center}
	
	\vspace{0.4cm}

\begin{enumerate}
	
	\item[] \textbf{Problema 1}
	
	Dadas las siguientes proposiciones relacionadas con procesos cognitivos:
	
	
	\qquad $p:\ \text{La memoria a corto plazo tiene capacidad ilimitada.}$
	
	\qquad $q:\ \text{La memoria sensorial tiene una duración de varios minutos.}$
	
	\qquad $r:\ \text{La memoria a largo plazo es de duración prolongada.}$
	
	\qquad $s:\ \text{El individuo utiliza la información para tomar decisiones.}$
	
	\begin{enumerate}[label=\alph*)]
		
		\item \textbf{Traducción a lenguaje usual de la proposición compuesta}
		
		Se pide traducir:
		\[
		(p \to q)\land (\sim r \to s)
		\]
		
		\textbf{Paso 1:} Interpretamos cada parte.
		
		\[
		p \to q
		\]
		significa:
		
		\begin{quote}
			Si la memoria a corto plazo tiene capacidad ilimitada, entonces la memoria sensorial tiene una duración de varios minutos.
		\end{quote}
		
		\[
		\sim r \to s
		\]
		significa:
		
		\begin{quote}
			Si la memoria a largo plazo no es de duración prolongada, entonces el individuo utiliza la información para tomar decisiones.
		\end{quote}
		
		Como ambas proposiciones están unidas por el conector
		\[
		\land
		\]
		que significa ``y'', la traducción completa es:
		
		\begin{quote}
			\textbf{Si la memoria a corto plazo tiene capacidad ilimitada, entonces la memoria sensorial tiene una duración de varios minutos, y si la memoria a largo plazo no es de duración prolongada, entonces el individuo utiliza la información para tomar decisiones.}
		\end{quote}
		
		\textbf{Respuesta final:}
		\[
		(p \to q)\land (\sim r \to s)
		\]
		se traduce como:
		
		\begin{quote}
			Si la memoria a corto plazo tiene capacidad ilimitada, entonces la memoria sensorial tiene una duración de varios minutos, y si la memoria a largo plazo no es de duración prolongada, entonces el individuo utiliza la información para tomar decisiones.
		\end{quote}
		
		\vspace{0.4cm}
		
		\item \textbf{Determinación del valor de verdad de la proposición compuesta}
		
		Se debe analizar:
		\[
		[(p\land \sim q)\lor r]\to s
		\]
		
		\textbf{Paso 1:} Asignamos valores de verdad según el contexto.
		
		\begin{itemize}
			\item \(p\): ``La memoria a corto plazo tiene capacidad ilimitada''.
			
			Esto es \textbf{falso}, porque la memoria a corto plazo es limitada.
			
			\[
			p=F
			\]
			
			\item \(q\): ``La memoria sensorial tiene una duración de varios minutos''.
			
			Esto es \textbf{falso}, porque la memoria sensorial dura muy poco tiempo.
			
			\[
			q=F
			\]
			
			\item \(r\): ``La memoria a largo plazo es de duración prolongada''.
			
			Esto es \textbf{verdadero}.
			
			\[
			r=V
			\]
			
			\item \(s\): ``El individuo utiliza la información para tomar decisiones''.
			
			Esto se considera \textbf{verdadero}.
			
			\[
			s=V
			\]
		\end{itemize}
		
		\textbf{Paso 2:} Sustituimos en la expresión.
		
		\[
		[(p\land \sim q)\lor r]\to s
		\]
		
		\[
		=[(F\land \sim F)\lor V]\to V
		\]
		
		\textbf{Paso 3:} Calculamos la negación.
		
		\[
		\sim F = V
		\]
		
		Entonces:
		
		\[
		[(F\land V)\lor V]\to V
		\]
		
		\textbf{Paso 4:} Calculamos la conjunción.
		
		\[
		F\land V = F
		\]
		
		Entonces:
		
		\[
		[F\lor V]\to V
		\]
		
		\textbf{Paso 5:} Calculamos la disyunción.
		
		\[
		F\lor V = V
		\]
		
		Entonces:
		
		\[
		V\to V
		\]
		
		\textbf{Paso 6:} Evaluamos la implicación.
		
		Sabemos que:
		
		\[
		V\to V = V
		\]
		
		Por lo tanto:
		
		\[
		[(p\land \sim q)\lor r]\to s = V
		\]
		
		\textbf{Conclusión:}
		
		La proposición compuesta es \textbf{verdadera}.
		
		\textbf{Respuesta final:}
		\[
		[(p\land \sim q)\lor r]\to s
		\]
		tiene valor de verdad:
		\[
		\boxed{\text{Verdadero}}
		\]
		
	\end{enumerate}
	
\end{enumerate}

\end{document}